\Chapter{23}

\Verse{1}
{
    \zA{话说贾母次日仍领众人过节。}
}
{
    \eA{Soon after the day on which Chia Yuan-ch'un honoured the garden of Broad
        Vista with a visit, and her return to the Palace,}
}
{
}

\Verse{2}
{
    \zA{那元妃却自幸大观园回宫去后，便命将那日所有 的题咏，命探春抄录妥协，自己编次优劣，又令在大观园勒石，为千古风流雅事。}
}
{
    \eA{so our story goes, she forthwith desired that T'an-ch'un should make a
        careful copy, in consecutive order, of the verses, which had been composed
        and read out on that occasion, in order that she herself should assign them
        their rank, and adjudge the good and bad. And she also directed that an
        inscription should be engraved on a stone,}
}
{
}

\Verse{3}
{
    \zA{因此贾政命人选拔精工，大观园磨石镌字。}
}
{
    \eA{in the Broad Vista park, to serve in future years as a record of the
        pleasant and felicitous event; and Chia Cheng, therefore,}
}
{
}

\Verse{4}
{
    \zA{贾珍率领贾蓉贾蔷等监工。}
}
{
    \eA{gave orders to servants to go far and wide, and select skilful artificers
        and renowned workmen,}
}
{
}

\Verse{5}
{
    \zA{因贾蔷又管 着文官等十二个女戏子并行头等事，不得空闲，因此又将贾菖、贾菱、贾萍唤来监 工。}
}
{
    \eA{to polish the stone and engrave the characters in the garden of Broad Vista;
        while Chia Chen put himself at the head of Chia Jung, Chia P'ing and others
        to superintend the work. And as Chia Se had, on the other hand, the control
        of Wen Kuan and the rest of the singing girls,}
}
{
}

\Verse{6}
{
    \zA{一日烫蜡钉朱，动起手来。}
}
{
    \eA{twelve in all, as well as of their costumes and other properties,}
}
{
}

\Verse{7}
{
    \zA{这也不在话下。}
}
{
    \eA{he had no leisure to attend to anything else,}
}
{
}

\Verse{8}
{
    \zA{且说那玉皇庙并达摩庵两处，一班的十二个小沙弥并十二个小道士，如今挪出 大观园来，贾政正想发到各庙去分住。}
}
{
    \eA{and consequently once again sent for Chia Ch'ang and Chia Ling to come and
        act as overseers. On a certain day, the works were taken in hand for rubbing
        the stones smooth with wax, for carving the inscription, and tracing it with
        vermilion, but without entering into details on these matters too minutely,
        we will return to the two places,}
}
{
}

\Verse{9}
{
    \zA{不想后街上住的贾芹之母杨氏，正打算到贾 政这边谋一个大小事件与儿子管管，也好弄些银钱使用，可巧听见这件事，便坐车 来求凤姐。}
}
{
    \eA{the Yu Huang temple and the Ta Mo monastery. The company of twelve young
        bonzes and twelve young Taoist priests had now moved out of the Garden of
        Broad Vista, and Chia Cheng was meditating upon distributing them to various
        temples to live apart, when unexpectedly Chia Ch'in's mother, née Chou,--who
        resided in the back street, and had been at the time contemplating to pay a
        visit to Chia Cheng on this side so as to obtain some charge,}
}
{
}

\Verse{10}
{
    \zA{凤姐因见他素日嘴头儿乖滑，便依允了。}
}
{
    \eA{be it either large or small, for her son to look after, that he too should
        be put in the way of turning up some money to meet his expenses with,--came,}
}
{
}

\Verse{11}
{
    \zA{想了几句话，便回了王夫人说： “这些小和尚小道士万不可打发到别处去，一时娘娘出来，就要应承的。}
}
{
    \eA{as luck would have it, to hear that some work was in hand in this mansion,
        and lost no time in driving over in a curricle and making her appeal to lady
        Feng. And as lady Feng remembered that she had all along not presumed on her
        position to put on airs, she willingly acceded to her request, and after
        calling to memory some suitable remarks,}
}
{
}

\Verse{12}
{
    \zA{倘或散了， 若再用时，可又费事。}
}
{
    \eA{she at once went to make her report to madame Wang: "These young bonzes and
        Taoist priests,"}
}
{
}

\Verse{13}
{
    \zA{依我的主意，不如将他们都送到家庙铁槛寺去，月间不过派 一个人拿几两银子去买柴米就是了。}
}
{
    \eA{she said, "can by no means be sent over to other places; for were the
        Imperial consort to come out at an unexpected moment, they would then be
        required to perform services; and in the event of their being scattered,
        there will, when the time comes to requisition their help,}
}
{
}

\Verse{14}
{
    \zA{说声用，走去叫一声就来，一点儿不费事。”}
}
{
    \eA{again be difficulties in the way; and my idea is that it would be better to
        send them all to the family temple, the Iron Fence Temple; and every month
        all there will be to do will be to depute some one to take over a few taels
        for them to buy firewood and rice with,}
}
{
}

\Verse{15}
{
    \zA{王夫人听了，便商之于贾政。}
}
{
    \eA{that's all, and when there's even a sound of their being required uttered,}
}
{
}

\Verse{16}
{
    \zA{贾政听了笑道：“倒是提醒了我。}
}
{
    \eA{some one can at once go and tell them just one word 'come,' and they will
        come without the least trouble!"}
}
{
}

\Verse{17}
{
    \zA{就是这样。”}
}
{
    \eA{Madame Wang gave a patient ear to this proposal,}
}
{
}

\Verse{18}
{
    \zA{即时 唤贾琏。}
}
{
    \eA{and, in due course, consulted with Chia Cheng.}
}
{
}

\Verse{19}
{
    \zA{贾琏正同凤姐吃饭，一闻呼唤，放下饭便走。}
}
{
    \eA{"You've really," smiled Chia Cheng at these words, "reminded me how I should
        act! Yes, let this be done!" And there and then he sent for Chia Lien.}
}
{
}

\Verse{20}
{
    \zA{凤姐一把拉住，笑道：“你 先站住，听我说话：要是别的事，我不管；要是为小和尚小道士们的事，好歹你依 着我这么着。”}
}
{
    \eA{Chia Lien was, at the time, having his meal with lady Feng, but as soon as
        he heard that he was wanted, he put by his rice and was just walking off,
        when lady Feng clutched him and pulled him back. "Wait a while," she
        observed with a smirk,}
}
{
}

\Verse{21}
{
    \zA{如此这般，教了一套话。}
}
{
    \eA{"and listen to what I've got to tell you! if it's about anything else, I've
        nothing to do with it;}
}
{
}

\Verse{22}
{
    \zA{贾琏摇头笑道：“我不管!你有本事你说 去。”}
}
{
    \eA{but if it be about the young bonzes and young Taoists, you must, in this
        particular matter, please comply with this suggestion of mine,"}
}
{
}

\Verse{23}
{
    \zA{凤姐听说，把头一梗，把筷子一放，腮上带笑不笑的瞅着贾琏道：“你是真 话，还是玩话儿？”}
}
{
    \eA{after which, she went on in this way and that way to put him up to a whole
        lot of hints. "I know nothing about it," Chia Lien rejoined smilingly, "and
        as you have the knack you yourself had better go and tell him!" But as soon
        as lady Feng heard this remark,}
}
{
}

\Verse{24}
{
    \zA{贾琏笑道：“西廊下五嫂子的儿子芸儿求了我两三遭，要件事 管管，我应了，叫他等着。}
}
{
    \eA{she stiffened her head and threw down the chopsticks; and, with an
        expression on her cheeks, which looked like a smile and yet not a smile, she
        glanced angrily at Chia Lien. "Are you speaking in earnest," she inquired,
        "or are you only jesting?"}
}
{
}

\Verse{25}
{
    \zA{好容易出来这件事，你又夺了去！”}
}
{
    \eA{"Yün Erh, the son of our fifth sister-in-law of the western porch, has come
        and appealed to me two or three times,}
}
{
}

\Verse{26}
{
    \zA{凤姐儿笑道：“你 放心。}
}
{
    \eA{asking for something to look after," Chia Lien laughed, "and I assented and
        bade him wait;}
}
{
}

\Verse{27}
{
    \zA{园子东北角上，娘娘说了，还叫多多的种松柏树，楼底下还叫种些花草儿。}
}
{
    \eA{and now, after a great deal of trouble, this job has turned up; and there
        you are once again snatching it away!" "Compose your mind," lady Feng
        observed grinning, "for the Imperial Consort has hinted that directions
        should be given for the planting,}
}
{
}

\Verse{28}
{
    \zA{等这件事出来，我包管叫芸儿管这工程就是了。”}
}
{
    \eA{in the north-east corner of the park, of a further plentiful supply of pine
        and cedar trees, and that orders should also be issued for the addition,
        round the base of the tower,}
}
{
}

\Verse{29}
{
    \zA{贾琏道：“这也罢了。”}
}
{
    \eA{of a large number of flowers and plants and such like; and when this job
        turns up,}
}
{
}

\Verse{30}
{
    \zA{因又悄 悄的笑道：“我问你，我昨儿晚上不过要改个样儿，你为什么就那么扭手扭脚的 呢？”}
}
{
    \eA{I can safely tell you that Yun Erh will be called to assume control of these
        works." "Well if that be really so," Chia Lien rejoined, "it will after all
        do! But there's only one thing; all I was up to last night was simply to
        have some fun with you, but you obstinately and perversely wouldn't."}
}
{
}

\Verse{31}
{
    \zA{凤姐听了，把脸飞红，“嗤”的一笑，向贾琏啐了一口，依旧低下头吃饭。}
}
{
    \eA{Lady Feng, upon hearing these words, burst out laughing with a sound of
        Ch'ih, and spurting disdainfully at Chia Lien, she lowered her head and went
        on at once with her meal; during which time Chia Lien speedily walked away
        laughing the while,}
}
{
}

\Verse{32}
{
    \zA{贾琏笑着一径去了。}
}
{
    \eA{and betook himself to the front, where he saw Chia Cheng.}
}
{
}

\Verse{33}
{
    \zA{走到前面见了贾政，果然为小和尚的事。}
}
{
    \eA{It was, indeed, about the young bonzes, and Chia Lien readily carried out
        lady Feng's suggestion. "As from all appearances," he continued, "Ch'in Erh
        has,}
}
{
}

\Verse{34}
{
    \zA{贾琏便依着凤姐的话，说道：“看来 芹儿倒出息了，这件事竟交给他去管，横竖照里头的规例，每月支领就是了。”}
}
{
    \eA{actually, so vastly improved, this job should, after all, be entrusted to
        his care and management; and provided that in observance with the inside
        custom Ch'in Erh were each day told to receive the advances, things will go
        on all right." And as Chia Cheng had never had much attention to give to
        such matters of detail,}
}
{
}

\Verse{35}
{
    \zA{贾 政原不大理论这些小事，听贾琏如此说，便依允了。}
}
{
    \eA{he, as soon as he heard what Chia Lien had to say, immediately signified his
        approval and assent. And Chia Lien, on his return to his quarters,
        communicated the issue to lady Feng;}
}
{
}

\Verse{36}
{
    \zA{贾琏回房告诉凤姐，凤姐即命 人去告诉杨氏，贾芹便来见贾琏夫妻，感谢不尽。}
}
{
    \eA{whereupon lady Feng at once sent some one to go and notify dame Chou. Chia
        Ch'in came, in due course, to pay a visit to Chia Lien and his wife, and was
        incessant in his expressions of gratitude; and lady Feng bestowed upon him a
        further favour by giving him,}
}
{
}

\Verse{37}
{
    \zA{凤姐又做情先支三个月的费用， 叫他写了领字，贾琏画了押，登时发了对牌出去，银库上按数发出三个月的供给来， 白花花三百两。}
}
{
    \eA{as a first instalment, an advance of the funds necessary for three months'
        outlay, for which she bade him write a receipt; while Chia Lien filled up a
        cheque and signed it; and a counter-order was simultaneously issued, and he
        came out into the treasury where the sum specified for three months'
        supplies, amounting to three hundred taels,}
}
{
}

\Verse{38}
{
    \zA{贾芹随手拈了一块与掌平的人，叫他们“喝了茶罢”。}
}
{
    \eA{was paid out in pure ingots. Chia Ch'in took the first piece of silver that
        came under his hand, and gave it to the men in charge of the scales, with
        which he told them to have a cup of tea,}
}
{
}

\Verse{39}
{
    \zA{于是命小厮 拿了回家，与母亲商议。}
}
{
    \eA{and bidding, shortly after, a boy-servant take the money to his home, he
        held consultation with his mother;}
}
{
}

\Verse{40}
{
    \zA{登时雇车坐上，又雇了几辆车子至荣国府角门前，唤出二 十四个人来，坐上车子，一径往城外铁槛寺去了。}
}
{
    \eA{after which, he hired a donkey for himself to ride on, and also bespoke
        several carriages, and came to the back gate of the Jung Kuo mansion; where
        having called out the twenty young priests, they got into the carriages, and
        sped straightway beyond the city walls,}
}
{
}

\Verse{41}
{
    \zA{当下无话。}
}
{
    \eA{to the Temple of the Iron Fence,}
}
{
}

\Verse{42}
{
    \zA{如今且说那元妃在宫中编次《大观园题咏》，忽然想起那园中的景致，自从幸 过之后，贾政必定敬谨封锁，不叫人进去，岂不辜负此园？}
}
{
    \eA{where nothing of any note transpired at the time. But we will now notice
        Chia Yüan-ch'un, within the precincts of the Palace. When she had arranged
        the verses composed in the park of Broad Vista in their order of merit, she
        suddenly recollected that the sights in the garden were sure, ever since her
        visit through them, to be diligently and respectfully kept locked up by her
        father and mother;}
}
{
}

\Verse{43}
{
    \zA{况家中现有几个能诗会 赋的姊妹们，何不命他们进去居住，也不使佳人落魄，花柳无颜。}
}
{
    \eA{and that by not allowing any one to go in was not an injustice done to this
        garden? "Besides," (she pondered), "in that household, there are at present
        several young ladies, capable of composing odes, and able to write poetry,}
}
{
}

\Verse{44}
{
    \zA{却又想宝玉自幼 在姊妹丛中长大，不比别的兄弟，若不命他进去，又怕冷落了他，恐贾母王夫人心 上不喜，须得也命他进去居住方妥。}
}
{
    \eA{and why should not permission be extended to them to go and take their
        quarters in it; in order too that those winsome persons might not be
        deprived of good cheer, and that the flowers and willows may not lack any
        one to admire them!" But remembering likewise that Pao-yü had from his
        infancy grown up among that crowd of female cousins, and was such a contrast
        to the rest of his male cousins that were he not allowed to move into it,}
}
{
}

\Verse{45}
{
    \zA{命太监夏忠到荣府下一道谕：“命宝钗等在园 中居住，不可封锢；命宝玉也随进去读书。”}
}
{
    \eA{he would, she also apprehended, be made to feel forlorn; and dreading lest
        his grandmother and his mother should be displeased at heart, she thought it
        imperative that he too should be permitted to take up his quarters inside,
        so that things should be put on a satisfactory footing;}
}
{
}

\Verse{46}
{
    \zA{贾政王夫人接了谕命。}
}
{
    \eA{and directing the eunuch Hsia Chung to go to the Jung mansion and deliver
        her commands,}
}
{
}

\Verse{47}
{
    \zA{夏忠去后，便 回明贾母，遣人进去各处收拾打扫，安设帘幔床帐。}
}
{
    \eA{she expressed the wish that Pao-ch'ai and the other girls should live in the
        garden and that it should not be kept closed, and urged that Pao-yü should
        also shift into it, at his own pleasure, for the prosecution of his studies.}
}
{
}

\Verse{48}
{
    \zA{别人听了，还犹自可，惟宝玉喜之不胜。}
}
{
    \eA{And Chia Cheng and madame Wang, upon receiving her commands, hastened, after
        the departure of Hsia Chung, to explain them to dowager lady Chia,}
}
{
}

\Verse{49}
{
    \zA{正和贾母盘算要这个要那个，忽见丫 鬟来说：“老爷叫宝玉。”}
}
{
    \eA{and to despatch servants into the garden to tidy every place, to dust, to
        sweep, and to lay out the portieres and bed-curtains. The tidings were heard
        by the rest even with perfect equanimity, but Pao-yü was immoderately
        delighted;}
}
{
}

\Verse{50}
{
    \zA{宝玉呆了半晌，登时扫了兴，脸上转了色，便拉着贾母 扭的扭股儿糖似的，死也不敢去。}
}
{
    \eA{and he was engaged in deliberation with dowager lady Chia as to this
        necessary and to that requirement, when suddenly they descried a
        waiting-maid arrive, who announced: "Master wishes to see Pao-yü." Pao-yü
        gazed vacantly for a while. His spirits simultaneously were swept away;}
}
{
}

\Verse{51}
{
    \zA{贾母只得安慰他道：“好宝贝，你只管去，有我 呢。}
}
{
    \eA{his countenance changed colour; and clinging to old lady Chia, he readily
        wriggled her about, just as one would twist the sugar (to make sweetmeats
        with),}
}
{
}

\Verse{52}
{
    \zA{他不敢委屈了你。}
}
{
    \eA{and could not, for the very death of him, summon up courage to go;}
}
{
}

\Verse{53}
{
    \zA{况你做了这篇好文章，想必娘娘叫你进园去住，他吩咐你几 句话，不过是怕你在里头淘气。}
}
{
    \eA{so that her ladyship had no alternative but to try and reassure him. "My
        precious darling" she urged, "just you go, and I'll stand by you! He won't
        venture to be hard upon you; and besides, you've devised these excellent
        literary compositions;}
}
{
}

\Verse{54}
{
    \zA{他说什么，你只好生答应着就是了。”}
}
{
    \eA{and I presume as Her Majesty has desired that you should move into the
        garden, his object is to give you a few words of advice;}
}
{
}

\Verse{55}
{
    \zA{一面安慰， 一面唤了两个老嬷嬷来，吩咐：“好生带了宝玉去，别叫他老子唬着他。”}
}
{
    \eA{simply because he fears that you might be up to pranks in those grounds. But
        to all he tells you, whatever you do, mind you acquiesce and it will be all
        right!" And as she tried to compose him, she at the same time called two old
        nurses and enjoined them to take Pao-yü over with due care,}
}
{
}

\Verse{56}
{
    \zA{老嬷嬷 答应了。}
}
{
    \eA{"And don't let his father," she added, "frighten him!"}
}
{
}

\Verse{57}
{
    \zA{宝玉只得前去，一步挪不了三寸，蹭到这边来。}
}
{
    \eA{The old nurses expressed their obedience, and Pao-yü felt constrained to
        walk ahead; and with one step scarcely progressing three inches,}
}
{
}

\Verse{58}
{
    \zA{可巧贾政在王夫人房中商议事情，金钏儿、彩云、彩凤、绣鸾、绣凤等众丫鬟 都廊檐下站着呢，一见宝玉来，都抿着嘴儿笑他。}
}
{
    \eA{he leisurely came over to this side. Strange coincidence Chia Cheng was in
        madame Wang's apartments consulting with her upon some matter or other, and
        Chin Ch'uan-erh, Ts'ai Yun, Ts'ai Feng, Ts'ai Luan, Hsiu Feng and the whole
        number of waiting-maids were all standing outside under the verandah. As
        soon as they caught sight of Pao-yü,}
}
{
}

\Verse{59}
{
    \zA{金钏儿一把拉着宝玉，悄悄的说 道：“我这嘴上是才擦的香香甜甜的胭脂，你这会子可吃不吃了？”}
}
{
    \eA{they puckered up their mouths and laughed at him; while Chin Ch'uan grasped
        Pao-yü with one hand, and remarked in a low tone of voice: "On these lips of
        mine has just been rubbed cosmetic, soaked with perfume, and are you now
        inclined to lick it or not?"}
}
{
}

\Verse{60}
{
    \zA{彩云一把推开 金钏儿，笑道：“人家心里发虚，你还怄他!趁这会子喜欢，快进去罢。”}
}
{
    \eA{whereupon Ts'ai Yün pushed off Chin Ch'uan with one shove, as she interposed
        laughingly, "A person's heart is at this moment in low spirits and do you
        still go on cracking jokes at him? But avail yourself of this opportunity
        when master is in good cheer to make haste and get in!"}
}
{
}

\Verse{61}
{
    \zA{宝玉只 得挨门进去。}
}
{
    \eA{Pao-yü had no help but to sidle against the door and walk in.}
}
{
}

\Verse{62}
{
    \zA{原来贾政和王夫人都在里间呢。}
}
{
    \eA{Chia Cheng and madame Wang were, in fact, both in the inner rooms, and dame
        Chou raised the portière.}
}
{
}

\Verse{63}
{
    \zA{赵姨娘打起帘子来，宝玉挨身而入， 只见贾政和王夫人对坐在炕上说话儿，地下一溜椅子，迎春、探春、惜春、贾环四 人都坐在那里。}
}
{
    \eA{Pao-yü stepped in gingerly and perceived Chia Cheng and madame Wang sitting
        opposite to each other, on the stove-couch, engaged in conversation; while
        below on a row of chairs sat Ying Ch'un, T'an Ch'un, Hsi Ch'un and Chia
        Huan; but though all four of them were seated in there only T'an Ch'un, Hsi
        Ch'un and Chia Huan rose to their feet,}
}
{
}

\Verse{64}
{
    \zA{一见他进来，探春惜春和贾环都站起来。}
}
{
    \eA{as soon as they saw him make his appearance in the room; and when Chia Cheng
        raised his eyes and noticed Pao-yü standing in front of him,}
}
{
}

\Verse{65}
{
    \zA{贾政一举目见宝玉站在跟前，神彩飘逸，秀色夺人，又看看贾环人物委琐，举 止粗糙，忽又想起贾珠来。}
}
{
    \eA{with a gait full of ease and with those winsome looks of his, so
        captivating, he once again realised what a mean being Chia Huan was, and how
        coarse his deportment. But suddenly he also bethought himself of Chia Chu,
        and as he reflected too that madame Wang had only this son of her own flesh
        and blood,}
}
{
}

\Verse{66}
{
    \zA{再看看王夫人只有这一个亲生的儿子，素爱如珍；自己 的胡须将已苍白：因此上把平日嫌恶宝玉之心不觉减了八九分。}
}
{
    \eA{upon whom she ever doated as upon a gem, and that his own beard had already
        begun to get hoary, the consequence was that he unwittingly stifled, well
        nigh entirely, the feeling of hatred and dislike, which, during the few
        recent years he had ordinarily fostered towards Pao-yü.}
}
{
}

\Verse{67}
{
    \zA{半晌说道：“娘娘 吩咐说：你日日在外游嬉，渐次疏懒了工课，如今叫禁管你和姐妹们在园里读书。}
}
{
    \eA{And after a long pause, "Her Majesty," he observed, "bade you day after day
        ramble about outside to disport yourself, with the result that you gradually
        became remiss and lazy; but now her desire is that we should keep you under
        strict control,}
}
{
}

\Verse{68}
{
    \zA{你可好生用心学习，再不守分安常，你可仔细着！”}
}
{
    \eA{and that in prosecuting your studies in the company of your cousins in the
        garden, you should carefully exert your brains to learn;}
}
{
}

\Verse{69}
{
    \zA{宝玉连连答应了几个“是”。}
}
{
    \eA{so that if you don't again attend to your duties, and mind your regular
        tasks,}
}
{
}

\Verse{70}
{
    \zA{王夫人便拉他在身边坐下。}
}
{
    \eA{you had better be on your guard!" Pao-yü assented several consecutive yes's;}
}
{
}

\Verse{71}
{
    \zA{他姊弟三人依旧坐下，王夫人摸索着宝玉的脖项说道： “前儿的丸药都吃完了没有？”}
}
{
    \eA{whereupon madame Wang drew him by her side and made him sit down, and while
        his three cousins resumed the seats they previously occupied: "Have you
        finished all the pills you had been taking a short while back?" madame Wang
        inquired, as she rubbed Pao-yü's neck.}
}
{
}

\Verse{72}
{
    \zA{宝玉答应道：“还有一丸。”}
}
{
    \eA{"There's still one pill remaining," Pao-yü explained by way of reply.}
}
{
}

\Verse{73}
{
    \zA{王夫人道：“明儿再 取十丸来，天天临睡时候，叫袭人伏侍你吃了再睡。”}
}
{
    \eA{"You had better," madame Wang added, "fetch ten more pills tomorrow morning;
        and every day about bedtime tell Hsi Jen to give them to you; and when
        you've had one you can go to sleep!"}
}
{
}

\Verse{74}
{
    \zA{宝玉道：“从太太吩咐了， 袭人天天临睡打发我吃的。”}
}
{
    \eA{"Ever since you, mother, bade me take them," Pao-yü rejoined, "Hsi Jen has
        daily sent me one, when I was about to turn in."}
}
{
}

\Verse{75}
{
    \zA{贾政便问道：“谁叫‘袭人’？”}
}
{
    \eA{"Who's this called Hsi Jen?" Chia Chen thereupon ascertained. "She's a
        waiting-maid!" madame Wang answered.}
}
{
}

\Verse{76}
{
    \zA{王夫人道：“是个 丫头。”}
}
{
    \eA{"A servant girl," Chia Cheng remonstrated, "can be called by whatever name
        one chooses; anything is good enough;}
}
{
}

\Verse{77}
{
    \zA{贾政道：“丫头不拘叫个什么罢了，是谁起这样刁钻名字？”}
}
{
    \eA{but who's it who has started this kind of pretentious name!" Madame Wang
        noticed that Chia Cheng was not in a happy frame of mind, so that she
        forthwith tried to screen matters for Pao-yü, by saying: "It's our old lady
        who has originated it!" "How can it possibly be," Chia Cheng exclaimed,
        "that her ladyship knows anything about such kind of language?}
}
{
}

\Verse{78}
{
    \zA{王夫人见贾 政不喜欢了，便替宝玉掩饰道：“是老太太起的。”}
}
{
    \eA{It must, for a certainty, be Pao-yü!" Pao-yü perceiving that he could not
        conceal the truth from him, was under the necessity of standing up and of
        explaining; "As I have all along read verses,}
}
{
}

\Verse{79}
{
    \zA{贾政道：“老太太如何晓得这 样的话？}
}
{
    \eA{I remembered the line written by an old poet: "What time the smell of
        flowers wafts itself into man, one knows the day is warm. "And as this
        waiting-maid's surname was Hua (flower),}
}
{
}

\Verse{80}
{
    \zA{一定是宝玉。”}
}
{
    \eA{I readily gave her the name, on the strength of this sentiment."}
}
{
}

\Verse{81}
{
    \zA{宝玉见瞒不过，只得起身回道：“因素日读诗，曾记古人 有句诗云：‘花气袭人知昼暖’，因这丫头姓‘花’，便随意起的。”}
}
{
    \eA{"When you get back," madame Wang speedily suggested addressing Pao-yü,
        "change it and have done; and you, sir, needn't lose your temper over such a
        trivial matter!" "It doesn't really matter in the least," Chia Cheng
        continued; "so that there's no necessity of changing it; but it's evident
        that Pao-yü doesn't apply his mind to legitimate pursuits,}
}
{
}

\Verse{82}
{
    \zA{王夫人忙向 宝玉说道：“你回去改了罢。}
}
{
    \eA{but mainly devotes his energies to such voluptuous expressions and wanton
        verses!" And as he finished these words, he abruptly shouted out:}
}
{
}

\Verse{83}
{
    \zA{老爷也不用为这小事生气。”}
}
{
    \eA{"You brute-like child of retribution! Don't you yet get out of this?" "Get
        away, off with you!"}
}
{
}

\Verse{84}
{
    \zA{贾政道：“其实也无妨 碍，不用改。}
}
{
    \eA{madame Wang in like manner hastened to urge; "our dowager lady is waiting, I
        fear, for you to have her repast!"}
}
{
}

\Verse{85}
{
    \zA{只可见宝玉不务正，专在这些浓词艳诗上做工夫。”}
}
{
    \eA{Pao-yü assented, and, with gentle step, he withdrew out of the room,
        laughing at Chin Ch'uan-erh, as he put out his tongue; and leading off the
        two nurses,}
}
{
}

\Verse{86}
{
    \zA{说毕，断喝了一 声：“作孽的畜生，还不出去！”}
}
{
    \eA{he went off on his way like a streak of smoke. But no sooner had he reached
        the door of the corridor than he espied Hsi Jen standing leaning against the
        side;}
}
{
}

\Verse{87}
{
    \zA{王夫人也忙道：“去罢，去罢。}
}
{
    \eA{who perceiving Pao-yü come back safe and sound heaped smile upon smile, and
        asked: "What did he want you for?" "There was nothing much,"}
}
{
}

\Verse{88}
{
    \zA{怕老太太等吃饭 呢。”}
}
{
    \eA{Pao-yü explained, "he simply feared that I would, when I get into the
        garden,}
}
{
}

\Verse{89}
{
    \zA{宝玉答应了，慢慢的退出去，向金钏儿笑着伸伸舌头，带着两个老嬷嬷，一溜 烟去了。}
}
{
    \eA{be up to mischief, and he gave me all sorts of advice;" and, as while he
        explained matters, they came into the presence of lady Chia, he gave her a
        clear account, from first to last, of what had transpired. But when he saw
        that Lin Tai-yü was at the moment in the room,}
}
{
}

\Verse{90}
{
    \zA{刚至穿堂门前，只见袭人倚门而立，见宝玉平安回来，堆下笑来，问道： “叫你做什么？”}
}
{
    \eA{Pao-yü speedily inquired of her: "Which place do you think best to live in?"
        Tai-yü had just been cogitating on this subject, so that when she
        unexpectedly heard Pao-yü's inquiry, she forthwith rejoined with a smile:
        "My own idea is that the Hsio Hsiang Kuan is best;}
}
{
}

\Verse{91}
{
    \zA{宝玉告诉：“没有什么，不过怕我进园淘气，吩咐吩咐。”}
}
{
    \eA{for I'm fond of those clusters of bamboos, which hide from view the tortuous
        balustrade and make the place more secluded and peaceful than any other!"
        Pao-yü at these words clapped his hands and smiled.}
}
{
}

\Verse{92}
{
    \zA{一面 说，一面回至贾母跟前回明原委。}
}
{
    \eA{"That just meets with my own views!" he remarked; "I too would like you to
        go and live in there; and as I am to stay in the I Hung Yuan,}
}
{
}

\Verse{93}
{
    \zA{只见黛玉正在那里，宝玉便问他：“你住在那一 处好？”}
}
{
    \eA{we two will be, in the first place, near each other; and next, both in quiet
        and secluded spots." While the two of them were conversing,}
}
{
}

\Verse{94}
{
    \zA{黛玉正盘算这事，忽见宝玉一问，便笑道：“我心里想着潇湘馆好。}
}
{
    \eA{a servant came, sent over by Chia Cheng, to report to dowager lady Chia
        that: "The 22nd of the second moon was a propitious day for Pao-yü and the
        young ladies to shift their quarters into the garden;}
}
{
}

\Verse{95}
{
    \zA{我爱 那几竿竹子，隐着一道曲栏，比别处幽静些。”}
}
{
    \eA{that during these few days, servants should be sent in to put things in
        their proper places and to clean; that Hsueh Pao-ch'ai should put up in the
        Heng Wu court;}
}
{
}

\Verse{96}
{
    \zA{宝玉听了，拍手笑道：“合了我的 主意了，我也要叫你那里住。}
}
{
    \eA{that Lin Tai-yü was to live in the Hsiao Hsiang lodge; that Chia Ying-ch'un
        should move into the Cho Chin two-storied building; that T'an Ch'un should
        put up in the Ch'iu Yen library;}
}
{
}

\Verse{97}
{
    \zA{我就住怡红院，咱们两个又近，又都清幽。”}
}
{
    \eA{that Hsi Ch'un should take up her quarters in the Liao Feng house; that
        widow Li should live in the Tao Hsiang village, and that Pao-yü was to live
        in the I Hung court.}
}
{
}

\Verse{98}
{
    \zA{二人正 计议着，贾政遣人来回贾母，说是：“二月二十二日是好日子，哥儿姐儿们就搬进 去罢。}
}
{
    \eA{That at every place two old nurses should be added and four servant-girls;
        that exclusive of the nurse and personal waiting-maid of each, there should,
        in addition, be servants, whose special duties should be to put things
        straight and to sweep the place; and that on the 22nd, they should all,}
}
{
}

\Verse{99}
{
    \zA{这几日便遣人进去分派收拾。”}
}
{
    \eA{in a body, move into the garden." When this season drew near, the interior
        of the grounds,}
}
{
}

\Verse{100}
{
    \zA{宝钗住了蘅芜院，黛玉住了潇湘馆，迎春住 了缀锦楼，探春住了秋掩书斋，惜春住了蓼风轩，李纨住了稻香村，宝玉住了怡红 院。}
}
{
    \eA{with the flowers waving like embroidered sashes, and the willows fanned by
        the fragrant breeze, was no more as desolate and silent as it had been in
        previous days; but without indulging in any further irrelevant details, we
        shall now go back to Pao-yü. Ever since he shifted his quarters into the
        park, his heart was full of joy,}
}
{
}

\Verse{101}
{
    \zA{每一处添两个老嬷嬷，四个丫头；除各人的奶娘亲随丫头外，另有专管收拾打 扫的。}
}
{
    \eA{and his mind of contentment, fostering none of those extraordinary ideas,
        whose tendency could be to give birth to longings and hankerings. Day after
        day, he simply indulged, in the company of his female cousins and the
        waiting-maids, in either reading his books,}
}
{
}

\Verse{102}
{
    \zA{至二十二日，一齐进去，登时园内花招绣带，柳拂香风，不似前番那等寂寞 了。}
}
{
    \eA{or writing characters, or in thrumming the lute, playing chess, drawing
        pictures and scanning verses, even in drawing patterns of argus pheasants,
        in embroidering phoenixes, contesting with them in searching for strange
        plants,}
}
{
}

\Verse{103}
{
    \zA{闲言少叙，且说宝玉自进园来，心满意足，再无别项可生贪求之心，每日只和 姊妹丫鬟们一处，或读书，或写字，或弹琴下棋，作画吟诗，以至描鸾刺凤，斗草
        簪花，低吟悄唱，拆字猜枚，无所不至，倒也十分快意。}
}
{
    \eA{and gathering flowers, in humming poetry with gentle tone, singing ballads
        with soft voice, dissecting characters, and in playing at mora, so that,
        being free to go everywhere and anywhere, he was of course completely happy.
        From his pen emanate four ballads on the times of the four seasons, which,
        although they could not be looked upon as first-rate,}
}
{
}

\Verse{104}
{
    \zA{他曾有几首四时即事诗， 虽不算好，却是真情真景。}
}
{
    \eA{afford anyhow a correct idea of his sentiments, and a true account of the
        scenery. The ballad on the spring night runs as follows: The silken
        curtains, thin as russet silk,}
}
{
}

\Verse{105}
{
    \zA{《春夜即事》云： 霞绡云幄任铺陈，隔巷蛙声听未真。}
}
{
    \eA{at random are spread out. The croak of frogs from the adjoining lane but
        faintly strikes the ear. The pillow a slight chill pervades, for rain
        outside the window falls.}
}
{
}

\Verse{106}
{
    \zA{枕上轻寒窗外雨，眼前春色梦中人。}
}
{
    \eA{The landscape, which now meets the eye, is like that seen in dreams by man.
        In plenteous streams the candles' tears do drop,}
}
{
}

\Verse{107}
{
    \zA{盈盈烛泪因谁泣，点点花愁为我嗔。}
}
{
    \eA{but for whom do they weep? Each particle of grief felt by the flowers is due
        to anger against me. It's all because the maids have by indulgence indolent
        been made.}
}
{
}

\Verse{108}
{
    \zA{自是小鬟妖懒惯，拥衾不耐笑言频。}
}
{
    \eA{The cover over me I'll pull, as I am loth to laugh and talk for long. This
        is the description of the aspect of nature on a summer night:}
}
{
}

\Verse{109}
{
    \zA{《夏夜即事》云： 倦绣佳人幽梦长，金笼鹦鹉唤茶汤。}
}
{
    \eA{The beauteous girl, weary of needlework, quiet is plunged in a long dream.
        The parrot in the golden cage doth shout that it is time the tea to brew.
        The lustrous windows with the musky moon like open palace-mirrors look;}
}
{
}

\Verse{110}
{
    \zA{窗明麝月开宫镜，室霭檀云品御香。}
}
{
    \eA{The room abounds with fumes of sandalwood and all kinds of imperial scents.
        From the cups made of amber is poured out the slippery dew from the lotus.}
}
{
}

\Verse{111}
{
    \zA{琥珀杯倾荷露滑，玻璃槛纳柳风凉。}
}
{
    \eA{The banisters of glass, the cool zephyr enjoy flapped by the willow trees.
        In the stream-spanning kiosk, the curtains everywhere all at one time do
        wave.}
}
{
}

\Verse{112}
{
    \zA{水亭处处齐纨动，帘卷朱楼罢晚妆。}
}
{
    \eA{In the vermilion tower the blinds the maidens roll, for they have made the
        night's toilette. The landscape of an autumnal evening is thus depicted:}
}
{
}

\Verse{113}
{
    \zA{《秋夜即事》云： 绛芸轩里绝喧哗，桂魄流光浸茜纱。}
}
{
    \eA{In the interior of the Chiang Yün house are hushed all clamorous din and
        noise. The sheen, which from Selene flows, pervades the windows of carnation
        gauze.}
}
{
}

\Verse{114}
{
    \zA{苔锁石纹容睡鹤，井飘桐露湿栖鸦。}
}
{
    \eA{The moss-locked, streaked rocks shelter afford to the cranes, plunged in
        sleep. The dew, blown on the t'ung tree by the well,}
}
{
}

\Verse{115}
{
    \zA{抱衾婢至舒金凤，倚槛人归落翠花。}
}
{
    \eA{doth wet the roosting rooks. Wrapped in a quilt, the maid comes the gold
        phoenix coverlet to spread. The girl, who on the rails did lean,}
}
{
}

\Verse{116}
{
    \zA{静夜不眠因酒渴，沉烟重拨索烹茶。}
}
{
    \eA{on her return drops the kingfisher flowers! This quiet night his eyes in
        sleep he cannot close, as he doth long for wine.}
}
{
}

\Verse{117}
{
    \zA{《冬夜即事》云： 梅魂竹梦已三更，锦？？}
}
{
    \eA{The smoke is stifled, and the fire restirred, when tea is ordered to be
        brewed. The picture of a winter night is in this strain:}
}
{
}

\Verse{118}
{
    \zA{衾睡未成。}
}
{
    \eA{The sleep of the plum trees,}
}
{
}

\Verse{119}
{
    \zA{松影一庭惟见鹤，梨花满地不闻莺。}
}
{
    \eA{the dream of the bamboos the third watch have already reached. Under the
        embroidered quilt and the kingfisher coverlet one can't sleep for the cold.
        The shadow of fir trees pervades the court, but cranes are all that meet the
        eye.}
}
{
}

\Verse{120}
{
    \zA{女奴翠袖诗怀冷，公子金貂酒力轻。}
}
{
    \eA{Both far and wide the pear blossom covers the ground, but yet the hawk
        cannot be heard. The wish, verses to write,}
}
{
}

\Verse{121}
{
    \zA{却喜侍儿知试茗，扫将新雪及时烹。}
}
{
    \eA{fostered by the damsel with the green sleeves, has waxéd cold. The master,
        with the gold sable pelisse, cannot endure much wine.}
}
{
}

\Verse{122}
{
    \zA{不说宝玉闲吟，且说这几首诗，当时有一等势利人，见是荣国府十二三岁的公 子做的，抄录出来，各处称颂；再有等轻薄子弟，爱上那风流妖艳之句，也写在扇
        头壁上，不时吟哦赏赞。}
}
{
    \eA{But yet he doth rejoice that his attendant knows the way to brew the tea.
        The newly-fallen snow is swept what time for tea the water must be boiled.
        But putting aside Pao-yü, as he leisurely was occupied in scanning some
        verses, we will now allude to all these ballads. There lived, at that time,
        a class of people, whose wont was to servilely court the influential and
        wealthy, and who, upon perceiving that the verses were composed by a young
        lad of the Jung Kuo mansion,}
}
{
}

\Verse{123}
{
    \zA{因此上竟有人来寻诗觅字，倩画求题。}
}
{
    \eA{of only twelve or thirteen years of age, had copies made, and taking them
        outside sang their praise far and wide.}
}
{
}

\Verse{124}
{
    \zA{这宝玉一发得意了， 每日家做这些外务。}
}
{
    \eA{There were besides another sort of light-headed young men, whose heart was
        so set upon licentious and seductive lines,}
}
{
}

\Verse{125}
{
    \zA{谁想静中生动，忽一日，不自在起来，这也不好，那也不好， 出来进去只是发闷。}
}
{
    \eA{that they even inscribed them on fans and screen-walls, and time and again
        kept on humming them and extolling them. And to the above reasons must
        therefore be ascribed the fact that persons came in search of stanzas and in
        quest of manuscripts,}
}
{
}

\Verse{126}
{
    \zA{园中那些女陔子，正是混沌世界天真烂熳之时，坐卧不避，嬉 笑无心，那里知宝玉此时的心事？}
}
{
    \eA{to apply for sketches and to beg for poetical compositions, to the
        increasing satisfaction of Pao-yü, who day after day, when at home, devoted
        his time and attention to these extraneous matters. But who would have
        anticipated that he could ever in his quiet seclusion have become a prey to
        a spirit of restlessness?}
}
{
}

\Verse{127}
{
    \zA{那宝玉不自在，便懒在园内，只想外头鬼混，却 痴痴的又说不出什么滋味来。}
}
{
    \eA{Of a sudden, one day he began to feel discontent, finding fault with this
        and turning up his nose at that; and going in and coming out he was simply
        full of ennui. And as all the girls in the garden were just in the prime of
        youth,}
}
{
}

\Verse{128}
{
    \zA{茗烟见他这样，因想与他开心，左思右想皆是宝玉玩 烦了的，只有一件，不曾见过。}
}
{
    \eA{and at a time of life when, artless and unaffected, they sat and reclined
        without regard to retirement, and disported themselves and joked without
        heed, how could they ever have come to read the secrets which at this time
        occupied a place in the heart of Pao-yü?}
}
{
}

\Verse{129}
{
    \zA{想毕便走到书坊内，把那古今小说，并那飞燕、合 德、则天、玉环的“外传”，与那传奇角本，买了许多，孝敬宝玉。}
}
{
    \eA{But so unhappy was Pao-yü within himself that he soon felt loth to stay in
        the garden, and took to gadding about outside like an evil spirit; but he
        behaved also the while in an idiotic manner. Ming Yen, upon seeing him go on
        in this way, felt prompted, with the idea of affording his mind some
        distraction, to think of this and to devise that expedient;}
}
{
}

\Verse{130}
{
    \zA{宝玉一看，如 得珍宝。}
}
{
    \eA{but everything had been indulged in with surfeit by Pao-yü, and there was
        only this resource,}
}
{
}

\Verse{131}
{
    \zA{茗烟又嘱咐道：“不可拿进园去，叫人知道了，我就‘吃不了兜着走’了。”}
}
{
    \eA{(that suggested itself to him,) of which Pao-yü had not as yet had any
        experience. Bringing his reflections to a close, he forthwith came over to a
        bookshop, and selecting novels, both of old and of the present age,}
}
{
}

\Verse{132}
{
    \zA{宝玉那里肯不拿进去？}
}
{
    \eA{traditions intended for outside circulation on Fei Yen, Ho Te, Wu Tse-t'ien,
        and Yang Kuei-fei,}
}
{
}

\Verse{133}
{
    \zA{踟蹰再四，单把那文理雅道些的，拣了几套进去，放在床顶 上，无人时方看；那粗俗过露的，都藏于外面书房内。}
}
{
    \eA{as well as books of light literature consisting of strange legends, he
        purchased a good number of them with the express purpose of enticing Pao-yü
        to read them. As soon as Pao-yü caught sight of them, he felt as if he had
        obtained some gem or jewel. "But you mustn't," Ming Yen went on to enjoin
        him, "take them into the garden;}
}
{
}

\Verse{134}
{
    \zA{那日正当三月中浣，早饭后，宝玉携了一套《会真记》，走到沁芳闸桥那边桃 花底下一块石上坐着，展开《会真记》，从头细看。}
}
{
    \eA{for if any one were to come to know anything about them, I shall then suffer
        more than I can bear; and you should, when you go along, hide them in your
        clothes!" But would Pao-yü agree to not introducing them into the garden? So
        after much wavering, he picked out only several volumes of those whose style
        was more refined,}
}
{
}

\Verse{135}
{
    \zA{正看到“落红成阵”，只见一 阵风过，树上桃花吹下一大斗来，落得满身满书满地皆是花片。}
}
{
    \eA{and took them in, and threw them over the top of his bed for him to peruse
        when no one was present; while those coarse and very indecent ones, he
        concealed in a bundle in the outer library. On one day, which happened to be
        the middle decade of the third moon,}
}
{
}

\Verse{136}
{
    \zA{宝玉要抖将不来， 恐怕脚步践踏了，只得兜了那花瓣儿，来至池边，抖在池内。}
}
{
    \eA{Pao-yü, after breakfast, took a book, the "Hui Chen Chi," in his hand and
        walked as far as the bridge of the Hsin Fang lock. Seating himself on a
        block of rock, that lay under the peach trees in that quarter,}
}
{
}

\Verse{137}
{
    \zA{那花瓣儿浮在水面， 飘飘荡荡，竟流出沁芳闸去了。}
}
{
    \eA{he opened the Hui Chen Chi and began to read it carefully from the
        beginning. But just as he came to the passage: "the falling red (flowers)
        have formed a heap,"}
}
{
}

\Verse{138}
{
    \zA{回来只见地下还有许多花瓣。}
}
{
    \eA{he felt a gust of wind blow through the trees, bringing down a whole bushel
        of peach blossoms;}
}
{
}

\Verse{139}
{
    \zA{宝玉正踟蹰间，只听 背后有人说道：“你在这里做什么？”}
}
{
    \eA{and, as they fell, his whole person, the entire surface of the book as well
        as a large extent of ground were simply bestrewn with petals of the
        blossoms. Pao-yü was bent upon shaking them down;}
}
{
}

\Verse{140}
{
    \zA{宝玉一回头，却是黛玉来了，肩上担着花锄， 花锄上挂着纱囊，手内拿着花帚。}
}
{
    \eA{but as he feared lest they should be trodden under foot, he felt constrained
        to carry the petals in his coat and walk to the bank of the pond and throw
        them into the stream. The petals floated on the surface of the water, and,
        after whirling and swaying here and there,}
}
{
}

\Verse{141}
{
    \zA{宝玉笑道：“来的正好，你把这些花瓣儿都扫起 来，撂在那水里去罢。}
}
{
    \eA{they at length ran out by the Hsin Fang lock. But, on his return under the
        tree, he found the ground again one mass of petals, and Pao-yü was just
        hesitating what to do, when he heard some one behind his back inquire,}
}
{
}

\Verse{142}
{
    \zA{我才撂了好些在那里了。”}
}
{
    \eA{"What are you up to here?" and as soon as Pao-yü turned his head round, he
        discovered that it was Lin Tai-yü,}
}
{
}

\Verse{143}
{
    \zA{黛玉道：“撂在水里不好，你看 这里的水干净，只一流出去，有人家的地方儿什么没有？}
}
{
    \eA{who had come over carrying on her shoulder a hoe for raking flowers, that on
        this hoe was suspended a gauze-bag, and that in her hand she held a broom.
        "That's right, well done!" Pao-yü remarked smiling; "come and sweep these
        flowers, and throw them into the water yonder.}
}
{
}

\Verse{144}
{
    \zA{仍旧把花遭塌了。}
}
{
    \eA{I've just thrown a lot in there myself!" "It isn't right," Lin Tai-yü
        rejoined,}
}
{
}

\Verse{145}
{
    \zA{那畸角 儿上我有一个花冢，如今把他扫了，装在这绢袋里，埋在那里；日久随土化了，岂 不干净。”}
}
{
    \eA{"to throw them into the water. The water, which you see, is clean enough
        here, but as soon as it finds its way out, where are situated other people's
        grounds, what isn't there in it? so that you would be misusing these flowers
        just as much as if you left them here! But in that corner, I have dug a hole
        for flowers,}
}
{
}

\Verse{146}
{
    \zA{宝玉听了，喜不自禁，笑道：“待我放下书，帮你来收拾。”}
}
{
    \eA{and I'll now sweep these and put them into this gauze-bag and bury them in
        there; and, in course of many days, they will also become converted into
        earth, and won't this be a clean way (of disposing of them)?"}
}
{
}

\Verse{147}
{
    \zA{黛玉道：“什么 书？”}
}
{
    \eA{Pao-yü, after listening to these words, felt inexpressibly delighted.}
}
{
}

\Verse{148}
{
    \zA{宝玉见问，慌的藏了，便说道：“不过是《中庸》《大学》。”}
}
{
    \eA{"Wait!" he smiled, "until I put down my book, and I'll help you to clear
        them up!" "What's the book?" Tai-yü inquired. Pao-yü at this question was so
        taken aback that he had no time to conceal it.}
}
{
}

\Verse{149}
{
    \zA{黛玉道：“你 又在我跟前弄鬼。}
}
{
    \eA{"It's," he replied hastily, "the Chung Yung and the Ta Hsüeh!" "Are you
        going again to play the fool with me?}
}
{
}

\Verse{150}
{
    \zA{趁早儿给我瞧瞧，好多着呢！”}
}
{
    \eA{Be quick and give it to me to see; and this will be ever so much better a
        way!" "Cousin," Pao-yü replied,}
}
{
}

\Verse{151}
{
    \zA{宝玉道：“妹妹，要论你我是不 怕的，你看了好歹别告诉人。}
}
{
    \eA{"as far as you yourself are concerned I don't mind you, but after you've
        seen it, please don't tell any one else. It's really written in beautiful
        style; and were you to once begin reading it,}
}
{
}

\Verse{152}
{
    \zA{真是好文章!你要看了，连饭也不想吃呢！”}
}
{
    \eA{why even for your very rice you wouldn't have a thought?" As he spoke, he
        handed it to her; and Tai-yü deposited all the flowers on the ground, took
        over the book, and read it from the very first page;}
}
{
}

\Verse{153}
{
    \zA{一面说， 一面递过去。}
}
{
    \eA{and the more she perused it, she got so much the more fascinated by it,}
}
{
}

\Verse{154}
{
    \zA{黛玉把花具放下，接书来瞧，从头看去，越看越爱，不顿饭时，已看 了好几出了。}
}
{
    \eA{that in no time she had finished reading sixteen whole chapters. But aroused
        as she was to a state of rapture by the diction, what remained even of the
        fascination was enough to overpower her senses; and though she had finished
        reading, she nevertheless continued in a state of abstraction,}
}
{
}

\Verse{155}
{
    \zA{但觉词句警人，馀香满口。}
}
{
    \eA{and still kept on gently recalling the text to mind, and humming it to
        herself. "Cousin, tell me is it nice or not?"}
}
{
}

\Verse{156}
{
    \zA{一面看了，只管出神，心内还默默记诵。}
}
{
    \eA{Pao-yü grinned. "It is indeed full of zest!" Lin Tai-yü replied exultingly.
        "I'm that very sad and very sickly person,"}
}
{
}

\Verse{157}
{
    \zA{宝玉笑道：“妹妹，你说好不好？”}
}
{
    \eA{Pao-yü explained laughing, "while you are that beauty who could subvert the
        empire and overthrow the city." Lin Tai-yü became, at these words,
        unconsciously crimson all over her cheeks,}
}
{
}

\Verse{158}
{
    \zA{黛玉笑着点头儿。}
}
{
    \eA{even up to her very ears; and raising, at the same moment,}
}
{
}

\Verse{159}
{
    \zA{宝玉笑道：“我就是个‘多 愁多病的身’，你就是那‘倾国倾城的貌’。”}
}
{
    \eA{her two eyebrows, which seemed to knit and yet not to knit, and opening wide
        those eyes, which seemed to stare and yet not to stare, while her peach-like
        cheeks bore an angry look and on her thin-skinned face lurked displeasure,
        she pointed at Pao-yü and exclaimed:}
}
{
}

\Verse{160}
{
    \zA{黛玉听了，不觉带腮连耳的通红了， 登时竖起两道似蹙非蹙的眉，瞪了一双似睁非睁的眼，桃腮带怒，薄面含嗔，指着
        宝玉道：“你这该死的，胡说了!好好儿的，把这些淫词艳曲弄了来，说这些混帐 话，欺负我。}
}
{
    \eA{"You do deserve death, for the rubbish you talk! without any provocation you
        bring up these licentious expressions and wanton ballads to give vent to all
        this insolent rot, in order to insult me; but I'll go and tell uncle and
        aunt." As soon as she pronounced the two words "insult me," her eyeballs at
        once were suffused with purple, and turning herself round she there and then
        walked away; which filled Pao-yü with so much distress that he jumped
        forward to impede her progress,}
}
{
}

\Verse{161}
{
    \zA{我告诉舅舅、舅母去！”}
}
{
    \eA{as he pleaded: "My dear cousin, I earnestly entreat you to spare me this
        time!}
}
{
}

\Verse{162}
{
    \zA{说到“欺负”二字，就把眼圈儿红了，转身 就走。}
}
{
    \eA{I've indeed said what I shouldn't; but if I had any intention to insult you,
        I'll throw myself to-morrow into the pond, and let the scabby-headed turtle
        eat me up,}
}
{
}

\Verse{163}
{
    \zA{宝玉急了，忙向前拦住道：“好妹妹，千万饶我这一遭儿罢!要有心欺负你， 明儿我掉在池子里，叫个癞头鼋吃了去，变个大忘八，等你明儿做了‘一品夫人’
        病老归西的时候儿，我往你坟上替你驼一辈子碑去。”}
}
{
    \eA{so that I become transformed into a large tortoise. And when you shall have
        by and by become the consort of an officer of the first degree, and you
        shall have fallen ill from old age and returned to the west, I'll come to
        your tomb and bear your stone tablet for ever on my back!" As he uttered
        these words, Lin Tai-yü burst out laughing with a sound of "pu ch'ih," and
        rubbing her eyes, she sneeringly remarked:}
}
{
}

\Verse{164}
{
    \zA{说的黛玉“扑嗤”的一声笑 了，一面揉着眼，一面笑道：“一般唬的这么个样儿，还只管胡说。}
}
{
    \eA{"I too can come out with this same tune; but will you now still go on
        talking nonsense? Pshaw! you're, in very truth, like a spear-head, (which
        looks) like silver, (but is really soft as) wax!" "Go on, go on!" Pao-yü
        smiled after this remark; "and what you've said, I too will go and tell!"}
}
{
}

\Verse{165}
{
    \zA{呸!原来也是 个‘银样蜡枪头’。”}
}
{
    \eA{"You maintain," Lin Tai-yü rejoined sarcastically, "that after glancing at
        anything you're able to recite it;}
}
{
}

\Verse{166}
{
    \zA{宝玉听了，笑道：“你说说，你这个呢？}
}
{
    \eA{and do you mean to say that I can't even do so much as take in ten lines
        with one gaze?" Pao-yü smiled and put his book away,}
}
{
}

\Verse{167}
{
    \zA{我也告诉去。”}
}
{
    \eA{urging: "Let's do what's right and proper,}
}
{
}

\Verse{168}
{
    \zA{黛 玉笑道：“你说你会‘过目成诵’，难道我就不能‘一目十行’了？”}
}
{
    \eA{and at once take the flowers and bury them; and don't let us allude to these
        things!" Forthwith the two of them gathered the fallen blossoms; but no
        sooner had they interred them properly than they espied Hsi Jen coming,}
}
{
}

\Verse{169}
{
    \zA{宝玉一面收 书，一面笑道：“正经快把花儿埋了罢，别提那些个了。”}
}
{
    \eA{who went on to observe: "Where haven't I looked for you? What! have you
        found your way as far as this! But our senior master, Mr. Chia She, over
        there isn't well; and the young ladies have all gone over to pay their
        respects,}
}
{
}

\Verse{170}
{
    \zA{二人便收拾落花。}
}
{
    \eA{and our old lady has asked that you should be sent over;}
}
{
}

\Verse{171}
{
    \zA{正才掩埋妥协，只见袭人走来，说道：“那里没找到？}
}
{
    \eA{so go back at once and change your clothes!" When Pao-yü heard what she
        said, he hastily picked up his books, and saying good bye to Tai-yü, he came
        along with Hsi Jen,}
}
{
}

\Verse{172}
{
    \zA{摸在这里来了!那边大老 爷身上不好，姑娘们都过去请安去了，老太太叫打发你去呢。}
}
{
    \eA{back into his room, where we will leave him to effect the necessary change
        in his costume. But during this while, Lin Tai-yü was, after having seen
        Pao-yü walk away, and heard that all her cousins were likewise not in their
        rooms,}
}
{
}

\Verse{173}
{
    \zA{快回去换衣裳罢。”}
}
{
    \eA{wending her way back alone, in a dull and dejected mood,}
}
{
}

\Verse{174}
{
    \zA{宝玉听了，忙拿了书，别了黛玉，同袭人回房换衣不提。}
}
{
    \eA{towards her apartment, when upon reaching the outside corner of the wall of
        the Pear Fragrance court, she caught, issuing from inside the walls, the
        harmonious strains of the fife and the melodious modulations of voices
        singing.}
}
{
}

\Verse{175}
{
    \zA{这里黛玉见宝玉去了，听见众姐妹也不在房中，自己闷闷的。}
}
{
    \eA{Lin Tai-yü readily knew that it was the twelve singing-girls rehearsing a
        play; and though she did not give her mind to go and listen, yet a couple of
        lines were of a sudden blown into her ears, and with such clearness,}
}
{
}

\Verse{176}
{
    \zA{正欲回房，刚走 到梨香院墙角外，只听见墙内笛韵悠扬，歌声婉转，黛玉便知是那十二个女孩子演 习戏文。}
}
{
    \eA{that even one word did not escape. Their burden was this: These troth are
        beauteous purple and fine carmine flowers, which in this way all round do
        bloom, And all together lie ensconced along the broken well, and the
        dilapidated wall! But the moment Lin Tai-yü heard these lines,}
}
{
}

\Verse{177}
{
    \zA{虽未留心去听，偶然两句吹到耳朵内，明明白白一字不落道：“原来是姹 紫嫣红开遍，似这般都付与断井颓垣。”}
}
{
    \eA{she was, in fact, so intensely affected and agitated that she at once halted
        and lending an ear listened attentively to what they went on to sing, which
        ran thus: A glorious day this is, and pretty scene, but sad I feel at heart!
        Contentment and pleasure are to be found in whose family courts? After
        overhearing these two lines,}
}
{
}

\Verse{178}
{
    \zA{黛玉听了，倒也十分感慨缠绵，便止步侧 耳细听。}
}
{
    \eA{she unconsciously nodded her head, and sighed, and mused in her own mind.
        "Really," she thought, "there is fine diction even in plays! but
        unfortunately what men in this world simply know is to see a play,}
}
{
}

\Verse{179}
{
    \zA{又唱道是：“良辰美景奈何天，赏心乐事谁家院。”}
}
{
    \eA{and they don't seem to be able to enjoy the beauties contained in them." At
        the conclusion of this train of thought, she experienced again a sting of
        regret,}
}
{
}

\Verse{180}
{
    \zA{听了这两句，不觉点 头自叹，心下自思：“原来戏上也有好文章，可惜世人只知看戏，未必能领略其中 的趣味。”}
}
{
    \eA{(as she fancied) she should not have given way to such idle thoughts and
        missed attending to the ballads; but when she once more came to listen, the
        song, by some coincidence, went on thus: It's all because thy loveliness is
        like a flower and like the comely spring, That years roll swiftly by just
        like a running stream.}
}
{
}

\Verse{181}
{
    \zA{想毕，又后悔不该胡想，耽误了听曲子。}
}
{
    \eA{When this couplet struck Tai-yu's ear, her heart felt suddenly a prey to
        excitement and her soul to emotion; and upon further hearing the words:}
}
{
}

\Verse{182}
{
    \zA{再听时，恰唱到：“只为你如 花美眷，似水流年。”}
}
{
    \eA{Alone you sit in the secluded inner rooms to self-compassion giving way.
        --and other such lines, she became still more as if inebriated, and like as
        if out of her head,}
}
{
}

\Verse{183}
{
    \zA{黛玉听了这两句，不觉心动神摇。}
}
{
    \eA{and unable to stand on her feet, she speedily stooped her body, and, taking
        a seat on a block of stone,}
}
{
}

\Verse{184}
{
    \zA{又听道“你在幽闺自怜” 等句，越发如醉如痴，站立不住，便一蹲身坐在一块山子石上，细嚼“如花美眷， 似水流年”八个字的滋味。}
}
{
    \eA{she minutely pondered over the rich beauty of the eight characters: It's all
        because thy loveliness is like a flower and like the comely spring, That
        years roll swiftly by just like a running stream. Of a sudden, she likewise
        bethought herself of the line: Water flows away and flowers decay, for both
        no feelings have. --which she had read some days back in a poem of an
        ancient writer, and also of the passage:}
}
{
}

\Verse{185}
{
    \zA{忽又想起前日见古人诗中，有“水流花谢两无情”之句； 再词中又有“流水落花春去也，天上人间”之句；又兼方才所见《西厢记》中“花
        落水流红，闲愁万种”之句：都一时想起来，凑聚在一处。}
}
{
    \eA{When on the running stream the flowers do fall, spring then is past and
        gone; --and of: Heaven (differs from) the human race, --which also appeared
        in that work; and besides these, the lines, which she had a short while back
        read in the Hsi Hiang Chi: The flowers, lo, fall, and on their course the
        waters red do flow! Petty misfortunes of ten thousand kinds (my heart
        assail!)}
}
{
}

\Verse{186}
{
    \zA{仔细忖度，不觉心痛神 驰，眼中落泪。}
}
{
    \eA{both simultaneously flashed through her memory; and, collating them all
        together, she meditated on them minutely,}
}
{
}

\Verse{187}
{
    \zA{正没个开交处，忽觉身背后有人拍了他一下，及至回头看时，未知是谁，下回 分解。}
}
{
    \eA{until suddenly her heart was stricken with pain and her soul fleeted away,
        while from her eyes trickled down drops of tears. But while nothing could
        dispel her present state of mind, she unexpectedly realised that some one
        from behind gave her a tap; and, turning her head round to look, she found
        that it was a young girl;}
}
{
}

\Verse{188}
{
    \zA{(本章完)}
}
{
    \eA{but who it was, the next chapter will make known.}
}
{
}
